\section{Risk Management and Mitigation}
\label{sec:risk}

\subsection{Attack surface}

Table~\ref{tab:surface} lists every entry point of the deployment in Figure~\ref{fig:arch}, who can
reach it, the threat it enables and the control we have. Digital surface dominates: there are no
user accounts and the physical surface is a single host.

\begin{table}[ht]
\centering\small
\caption{Attack surface of the k3s deployment.}
\label{tab:surface}
\begin{tabularx}{\textwidth}{@{}>{\raggedright}p{3.3cm} >{\raggedright}p{2.1cm} >{\raggedright}X >{\raggedright\arraybackslash}X@{}}
\toprule
Entry point & Reachable from & Threat & Control in place \\
\midrule
ingestion / analytics NodePorts & lab network & HTTP flood; forged readings & Wazuh 100210 + k3s-block; Prometheus load alerts; pod limits; input validation \\
registry API (no auth) & pod network only & fake or hijacked registrations (discovery poisoning) & ClusterIP only; health checks remove dead entries; \emph{no authentication} \\
storage API & pod network only & data tampering, read-heavy load & ClusterIP only; single writer \\
SMTP (mail server) & lab network & spam, mail flood, open relay & SpamAssassin + hold queue; 100222 + firewall-drop; relay restrictions \\
DNS zone file & DNS host & redirection of service names & real-time FIM + restore-dns \\
DNS resolution path & lab network & cache poisoning, rogue resolver & \emph{none} \\
k3s nodes (privileged containers) & host & container escape = host compromise & \emph{none}; documented limitation \\
Image registry (plain HTTP) & lab network & image tampering & lab only, \emph{no TLS or signing} \\
Prometheus UI/API (9090, NodePort 30090) & lab network & reading the topology; spotting the thresholds & \emph{none}: no authentication \\
Wazuh dashboard (443) & lab network & blinding or misleading the defender & login page; API (55000) and indexer (9200) not reachable \\
\bottomrule
\end{tabularx}
\end{table}

\subsection{Attack and defense tree}

Figure~\ref{fig:adtree} refines the attacker goal \emph{disrupt or subvert the temperature service}
into the attacks we considered, with our countermeasures attached as defence nodes, in the notation
of attack--defense trees~\cite{kordy2010adtrees}. All children are OR-refinements: each leaf alone
reaches its parent goal. Grey nodes mark attacks for which we have no defence; they are the
limitations discussed in Section~\ref{sec:concl}. The three threats of Section~\ref{sec:threats}
correspond to the DoS, DNS and junk-email branches.

\begin{figure}[ht]
\centering
\resizebox{\textwidth}{!}{% Figure 3: attack-defense tree (Kordy et al. notation: attacks = red ellipses, defences = green
% boxes on dotted edges; grey = no defence in our system). Every defence is an element of the repo
% (FACTS.md sections 4-7); every "no defence" leaf is a documented limitation (FACTS.md section 10).
\begin{forest}
  for tree={font=\scriptsize, align=center, grow'=east, l sep=7mm, s sep=1.5mm, inner sep=2pt,
            child anchor=west, parent anchor=east, anchor=west},
  att/.style={ellipse, draw=red!70!black, fill=red!8},
  dfn/.style={rectangle, draw=green!50!black, fill=green!10, edge={dotted, thick}},
  gap/.style={rectangle, draw=gray, fill=gray!15, edge={dotted, thick}, font=\scriptsize\itshape},
  [{Disrupt or subvert\\the temperature service}, att
    [{Deny service\\(DoS)}, att
      [{HTTP flood from\\few sources}, att
        [{100210 +\\k3s-block 60\,s}, dfn]
        [{Prometheus rate/\\latency alerts}, dfn]
      ]
      [{Distributed flood\\below 15 req/10\,s\\per source}, att
        [{no automatic\\response}, gap]
      ]
      [{Flood from a\\compromised pod}, att
        [{exempted by\\100209}, gap]
      ]
    ]
    [{Redirect clients\\(DNS hijack)}, att
      [{Tamper with\\zone file}, att
        [{FIM 550$\to$100230\\+ restore-dns}, dfn]
      ]
      [{Cache poisoning /\\rogue DNS}, att
        [{not detected\\(DNSSEC: future)}, gap]
      ]
    ]
    [{Poison service\\discovery}, att
      [{Register fake\\instance}, att
        [{registry only\\ClusterIP}, dfn]
        [{no authentication}, gap]
      ]
    ]
    [{Junk email}, att
      [{Spam content}, att
        [{SpamAssassin $\to$\\hold queue}, dfn]
      ]
      [{Mail flood\\one source}, att
        [{100222 +\\firewall-drop}, dfn]
      ]
      [{Open relay\\abuse}, att
        [{relay\\restrictions}, dfn]
      ]
    ]
    [{Take over\\the host}, att
      [{Escape the\\privileged k3s\\container}, att
        [{no control}, gap]
      ]
    ]
  ]
\end{forest}
}
\caption{Attack--defense tree. Red ellipses: attacks; green boxes: our defences; grey boxes:
no effective defence in the current system.}
\label{fig:adtree}
\end{figure}

The tree makes two priorities visible. First, the DoS defence covers only one of three ways to deny
service; the other two (many sources, a compromised pod) are only \emph{seen} by Prometheus.
Second, the registry is the most valuable unauthenticated asset: a fake registration redirects
traffic exactly like a DNS hijack, but is invisible to file-integrity monitoring.

\subsection{Large-scale monitoring strategy}

Prometheus scrapes every instance every 5~s. In the cluster, pods are discovered through the
Kubernetes API and relabelled to the same job names as on Incus (\code{ssle-registry},
\code{ssle-ingestion}, \dots), so one rule file serves both platforms. The rules have three groups:

\begin{itemize}
  \item \textbf{Load anomalies} (\code{ssle-ddos-alerts}, every 5~s): request rate above 5~req/s
        per job, p95 latency above 50~ms, and more than 3 requests in flight.
  \item \textbf{Correlation and root cause} (\code{ssle-root-cause}): each alert is labelled
        \code or \code{kind=root-cause} and names a \code{root\_cause}. For example,
        \code{SSLERootCauseStorageDown} fires only when ingestion or analytics report storage errors
        \emph{and} storage itself is down, while \code{SSLERootCauseStorageUnreachable} blames the
        network or discovery when storage is up but unreachable, unless a request flood explains the
        errors. \code{SSLERootCauseRequestFlood} combines high latency and high rate on the same job
        and points the operator to Wazuh rule \rid{100210} for the source addresses.
  \item \textbf{Per-location aggregates} (\code{ssle-location-aggregates}), used for federation
        (Section~\ref{sec:arch}).
\end{itemize}

Two rules were changed after the tests proved them wrong: during the DoS test the
\emph{unreachable} rule blamed the network for timeouts caused by an overloaded storage, so it now
requires the errors to persist and excludes floods; and on Kubernetes a pod scaled to zero has no
\code{up} series at all, so every rule that checks \code{up} uses \code{or vector(0)}.

\subsection{Fault management (FCAPS) and the security gap}

Figure~\ref{fig:fcaps} maps the FCAPS fault-management workflow~\cite{itu2000m3400} to our tools.
Prometheus covers the classic part: it \emph{detects} anomalies, the hierarchical ID
\emph{isolates} the element, the root-cause rules \emph{correlate} symptoms, the registry and
Kubernetes \emph{resolve} failed instances, and each test produces the \emph{post-mortem}. The
root-cause test of Section~\ref{sec:eval} exercises this chain end to end.

\begin{figure}[ht]
\centering
\resizebox{\textwidth}{!}{% Figure 4: FCAPS fault-management workflow mapped to our tooling (FACTS.md sections 5, 6, 8, 9).
\begin{tikzpicture}[
  font=\scriptsize,
  stage/.style={draw, rounded corners=3pt, fill=blue!12, minimum width=21mm, minimum height=7mm, align=center, font=\scriptsize\bfseries},
  tool/.style={draw, rounded corners=2pt, fill=gray!8, text width=21mm, align=left, font=\tiny},
  arr/.style={-{Stealth[length=2mm]}, thick, orange!70!black}
]
\node[stage] (d) at (0,0)     {Detecting};
\node[stage] (i) at (2.75,0)  {Diagnosing \&\\Isolating};
\node[stage] (c) at (5.5,0)   {Correlating};
\node[stage] (r) at (8.25,0)  {Resolution /\\Mitigation};
\node[stage] (p) at (11,0)    {Post-mortem};
\foreach \a/\b in {d/i,i/c,c/r,r/p} \draw[arr] (\a) -- (\b);

\node[tool, below=3mm of d] {Prometheus: \code{up}, rates, p95, in-flight (5\,s)\\Wazuh: access logs (100200), FIM (550), Postfix/spamd};
\node[tool, below=3mm of i] {hierarchical ID in every metric, alert and log line $\Rightarrow$ service, pod, node, rack, dc};
\node[tool, below=3mm of c] {\code{ssle-root-cause} rules labelled \code{kind=}\\\code{symptom} or \code{root-cause}\\Wazuh frequency rules 100210, 100222};
\node[tool, below=3mm of r] {active responses: k3s-block, firewall-drop, restore-dns\\registry drops dead instances; Kubernetes restarts pods};
\node[tool, below=3mm of p] {test scripts: threat card, timings, false positives, \code{summary.txt} per threat};

\node[font=\tiny\itshape, align=center, text width=60mm] at (5.5,0.85)
  {classic FCAPS fault management (Prometheus) \;+\; explicit security events (Wazuh)};
\end{tikzpicture}
}
\caption{FCAPS fault-management workflow and the tool that implements each step.}
\label{fig:fcaps}
\end{figure}

FCAPS fault management does not deal with explicit security events: from Prometheus' point of view a
DNS hijack is not a fault at all (every service is up and answering), and a flood is only "high
load". The SIEM closes this gap. Wazuh receives the same access logs plus file-integrity and mail
events, correlates them with frequency rules (\rid{100210}, \rid{100222}) and triggers active
responses. The two systems are linked by the common ID and by cross-references: the Prometheus flood
root cause points to the Wazuh rule, and Wazuh's HTTP~5xx events (\rid{100201}) from each service
are the security-side view of the same symptoms the root-cause rules correlate.
